\documentclass[11pt,a4paper]{article}

\usepackage[a4paper,left=0.85cm,right=0.85cm,top=0.85cm,bottom=0.75cm]{geometry}
\usepackage[T1]{fontenc}
\usepackage[utf8]{inputenc}
\usepackage{charter}
\usepackage[hidelinks]{hyperref}
\usepackage{enumitem}
\usepackage{tabularx}
\usepackage{xcolor}
\usepackage{titlesec}
\usepackage{array}

\definecolor{accent}{HTML}{123C55}
\definecolor{muted}{HTML}{555555}

\pagestyle{empty}
\setlength{\parindent}{0pt}
\setlength{\tabcolsep}{0pt}
\setlength{\parskip}{0pt}
\renewcommand{\arraystretch}{1.08}
\emergencystretch=1.5em

\titleformat{\section}
  {\large\bfseries\scshape\color{accent}}
  {}
  {0pt}
  {}
  [\vspace{0.10em}{\color{accent}\titlerule[0.55pt]}]
\titlespacing*{\section}{0pt}{0.72em}{0.42em}

\setlist[itemize]{
  leftmargin=1.15em,
  label=\raisebox{0.2ex}{\scriptsize$\bullet$},
  itemsep=0.06em,
  topsep=0.08em,
  parsep=0pt,
  partopsep=0pt
}

\newlength{\datecol}
\setlength{\datecol}{2.35cm}
\newlength{\gapcol}
\setlength{\gapcol}{0.35cm}
\newcommand{\entry}[3]{%
  \noindent
  \begin{minipage}[t]{\datecol}
    \vspace{0pt}\raggedright\small\color{muted}#1
  \end{minipage}%
  \hspace{\gapcol}%
  \begin{minipage}[t]{\dimexpr\linewidth-\datecol-\gapcol\relax}
    \vspace{0pt}\raggedright #2#3
  \end{minipage}\par\vspace{0.36em}
}

\newcommand{\role}[2]{%
  \textbf{#1}\if\relax\detokenize{#2}\relax\else, #2\fi\par
}

\newcommand{\project}[2]{%
  \noindent\textbf{#1}: #2\par\vspace{0.10em}
}

\begin{document}
\raggedright

{\LARGE\bfseries Rodrigo Alvarez Lucendo}\par
\vspace{0.25em}
{\small
\mbox{+41 77 289 46 00}
\quad|\quad
\mbox{\href{mailto:r.alvarezlucendo16@gmail.com}{r.alvarezlucendo16@gmail.com}}
\quad|\quad
\mbox{Spanish citizen}
}

\section{Education}
\entry{2023--2026}
{\role{MSc Computer Science}{TU Delft}}
{%
  \begin{itemize}
    \item Expected graduation: August 2026.
    \item Exchange semester at EPFL (Feb--Jul 2025).
    \item Courses: Modern NLP, Deep Learning, Graph ML, Deep RL, Computer Vision, Distributed ML Systems.
  \end{itemize}
}
\entry{2019--2023}
{\role{BSc Computer Science}{TU Delft}}
{%
  \begin{itemize}
    \item Thesis: Deep learning model for algal bloom forecasting using spatiotemporal data on HPC.
  \end{itemize}
}

\section{Work Experience}
\entry{Jun--Present}
{\role{Machine Learning Engineer}{Xcert.ai}}
{%
  \begin{itemize}
    \item First employee at CHF 1.5M+ startup building AI assistants for aerospace certification.
    \item Led RAG pipeline development (parsing, chunking, retrieval, and evaluation) as well as the development, 
    monitoring and evaluation of an agentic workflow created with OpenAI Agent SDK.
  \end{itemize}
}
\entry{Jul 2024--Jan 2025}
{\role{Data Engineer Intern}{Amazon}}
{%
  \begin{itemize}
    \item Deployed a monitoring tool for machine learning models in production.
    \item First in team to explore AI agents; built a ReAct agent for natural-language metric analysis.
    \item Provided BI teams with complex SQL queries and business analysis. 
  \end{itemize}
}
\entry{2022--2024}
{\role{Teaching Assistant}{TU Delft \& EPFL}}
{%
  \begin{itemize}
    \item Guided 300+ students in BSc/MSc courses (ML, Big Data, Deep Learning); Head TA for Computational Intelligence.
  \end{itemize}
}

\section{Research}
\entry{Feb--Sep 2025}
{\role{Research Assistant}{Theory of Machine Learning Lab, EPFL}}
{%
  \begin{itemize}
    \item Investigated how Transformers learn high-order Markov chains through staged dynamics.
  \end{itemize}
}
\entry{2025--2026}
{\role{Master Thesis}{Natural Language Processing Lab, EPFL}}
{%
  \begin{itemize}
    \item Measuring scaling laws for memorization in LLMs across model size and data composition.
    \item Training GPT-2 models from scratch on synthetic and natural data.
  \end{itemize}
}

\section{Projects}
\project{AI Tutor}{Fine-tuned a 0.6B LLM with SFT and DPO (Hugging Face) to create an AI STEM tutor.}
\project{nanoGPT}{Built nanoGPT from scratch (decoder-only Transformer, BPE tokenizer) and extended it with MoE.}
\project{RAG QA}{Analyzed how contextual relevance influences RAG QA system performance.}
\project{FETCH, EPFL AI Team}{Co-developed a natural-language mission planner for a Unitree Go2 robot using LLMs and live sensor input; integrated Nav2 and MCP servers.}

\section{Publications}
\textbf{Incremental Learning of Sparse Attention Patterns in Transformers.}
O\u{g}uz Kaan Y\"uksel, \textbf{Rodrigo Alvarez Lucendo}, Nicolas Flammarion.
Second author; accepted to ICML 2026 Main Conference (regular accept), attending as author.
Earlier version at EurIPS 2025 Workshop on Principles of Generative Modeling (PriGM).
\href{https://arxiv.org/abs/2602.19143}{arXiv:2602.19143}.

\section{Skills}
\begin{tabularx}{\linewidth}{@{}>{\bfseries}p{3.25cm}@{\hspace{0.35cm}}X@{}}
Languages & Python, Java, SQL, Scala, Bash, TypeScript, English, Spanish.\\
ML / Deep Tech & PyTorch, Hugging Face, LlamaIndex, PyTorch Lightning, Weights \& Biases.\\
Systems / Tools & Docker, AWS, Neo4j, Redis, FastAPI, Git, HPC.\\
Interests & Rowing, alpine skiing, strength training.
\end{tabularx}

\end{document}
